\documentclass[a4paper,12pt]{article}

% ---------- Минимальная стабильная преамбула ----------
\usepackage[T2A]{fontenc}
\usepackage[utf8]{inputenc}
\usepackage[russian]{babel}
\usepackage{amsmath,amssymb}
\usepackage{icomma}
\usepackage[a4paper,left=18mm,right=18mm,top=18mm,bottom=20mm]{geometry}
\usepackage{xcolor}
\usepackage{array,booktabs}
\usepackage{enumitem}
\usepackage{tikz}
\usepackage{pgfplots}
\pgfplotsset{compat=1.18}

% ---------- Оформление ----------
\renewcommand{\familydefault}{\sfdefault}
\setlength{\parindent}{0pt}
\setlength{\parskip}{5pt}
\linespread{1.12}
\emergencystretch=2em

\definecolor{accent}{HTML}{008080}
\definecolor{darkaccent}{HTML}{004D4D}
\definecolor{softgray}{HTML}{666666}
\definecolor{lightaccent}{HTML}{DDEEEE}

\newcommand{\key}[1]{\textcolor{darkaccent}{\textbf{#1}}}
\newcommand{\sepLine}{\par\vspace{2mm}{\color{lightaccent}\rule{\linewidth}{0.8pt}}\vspace{2mm}}

\setlist[itemize]{leftmargin=6mm,itemsep=2pt,topsep=2pt}
\setlist[enumerate]{leftmargin=8mm,itemsep=4pt,topsep=3pt}

% ---------- Единый стиль графиков ----------
\pgfplotsset{
  schoolaxis/.style={
    axis lines=middle,
    unit vector ratio=1 1,
    width=11.5cm,
    height=6.6cm,
    grid=major,
    xlabel={$x$},
    ylabel={$y$},
    samples=220,
    clip=false,
    xticklabels={},
    yticklabels={},
    tick style={draw=none},
    label style={font=\small}
  },
  maincurve/.style={accent,line width=1.2pt},
  secondcurve/.style={darkaccent,line width=1.2pt,dashed},
  knot/.style={only marks,mark=*,mark size=2.5pt,darkaccent}
}

\begin{document}

% ==================================================
% ТИТУЛЬНЫЙ ЛИСТ
% ==================================================
\begin{titlepage}
\thispagestyle{empty}
\begin{center}
\vspace*{30mm}

{\Large\textcolor{softgray}{\textbf{ЧЕК-ЛИСТ}}}

\vspace{10mm}

{\Huge\textbf{\textcolor{darkaccent}{Графики функций}}}

\vspace{5mm}

{\Large\textbf{Метод узловых точек}}

\vspace{5mm}

{\large Линейная функция, парабола, гипербола\par
и основные графики школьного курса}

\vspace{18mm}

{\color{lightaccent}\rule{0.65\linewidth}{1pt}}

\vspace{18mm}

{\large
\textbf{Лёвин Артём Александрович}\\[2mm]
эксклюзивно для Тимофея Васильченко}

\vfill

{\large \today}
\end{center}

\begin{tikzpicture}[remember picture,overlay]
\draw[accent,line width=1pt]
([xshift=18mm,yshift=15mm]current page.south west)
.. controls
([xshift=65mm,yshift=27mm]current page.south west)
and
([xshift=-65mm,yshift=6mm]current page.south east)
..
([xshift=-18mm,yshift=16mm]current page.south east);
\end{tikzpicture}
\end{titlepage}

% ==================================================
\section*{\textcolor{darkaccent}{1. Что нужно помнить}}

Вы должны уметь:

\begin{itemize}
  \item узнавать основные типы функций;
  \item восстанавливать формулу по узловым точкам;
  \item находить $f(a)$ и решать уравнение $f(x)=c$;
  \item находить координаты точек пересечения графиков;
  \item проверять ответ по смыслу и рисунку.
\end{itemize}

\sepLine

% ==================================================
\section*{\textcolor{darkaccent}{2. Основные понятия}}

Для точки
\[
A(x_A;y_A)
\]
координата $x_A$ называется \key{абсциссой}, а $y_A$ --- \key{ординатой}.

\[
\boxed{\text{абсцисса}=x},
\qquad
\boxed{\text{ордината}=y}.
\]

Запись
\[
f(a)
\]
означает: подставить $a$ вместо $x$.

Запись
\[
f(x)=c
\]
означает: решить уравнение относительно $x$.

\sepLine

% ==================================================
\section*{\textcolor{darkaccent}{3. Формулы основных функций}}

Формулы стандартных графиков полезно знать наизусть.

\begin{center}
\small
\renewcommand{\arraystretch}{1.25}
\begin{tabular}{@{}p{0.31\linewidth}p{0.31\linewidth}p{0.31\linewidth}@{}}
\toprule
\textbf{Функция} & \textbf{Формула} & \textbf{Главное} \\
\midrule
Постоянная & $y=c$ & горизонтальная прямая \\
Прямая пропорциональность & $y=kx$ & проходит через $(0;0)$ \\
Линейная & $y=kx+b$ & $k$ --- наклон, $b$ --- сдвиг \\
Квадратичная & $y=ax^2+bx+c$ & парабола, $a\ne0$ \\
Кубическая & $y=x^3$ & нечётная функция \\
Степенная & $y=x^n$ & вид зависит от $n$ \\
Обратная пропорциональность & $y=\dfrac{k}{x}$ & $x\ne0$ \\
Сдвинутая гипербола & $y=\dfrac{k}{x-a}+b$ & асимптоты $x=a$, $y=b$ \\
\bottomrule
\end{tabular}
\end{center}

\begin{center}
\small
\renewcommand{\arraystretch}{1.25}
\begin{tabular}{@{}p{0.31\linewidth}p{0.31\linewidth}p{0.31\linewidth}@{}}
\toprule
\textbf{Функция} & \textbf{Формула} & \textbf{Главное} \\
\midrule
Модуль & $y=|x|$ & V-образный график \\
Корень & $y=\sqrt{x}$ & $x\ge0$ \\
Экспонента & $y=a^x$ & $a>0$, $a\ne1$ \\
Логарифм & $y=\log_a x$ & $x>0$ \\
Синус & $y=\sin x$ & период $2\pi$ \\
Косинус & $y=\cos x$ & период $2\pi$ \\
Тангенс & $y=\tg x$ & период $\pi$ \\
\bottomrule
\end{tabular}
\end{center}

\subsection*{Сдвиги стандартных графиков}

Если известен график $y=f(x)$, то:

\[
f(x)+a
\quad\text{--- вертикальный перенос,}
\]

\[
f(x-a)
\quad\text{--- горизонтальный перенос,}
\]

\[
-f(x)
\quad\text{--- отражение относительно }Ox,
\]

\[
f(-x)
\quad\text{--- отражение относительно }Oy.
\]

\clearpage

% ==================================================
\section*{\textcolor{darkaccent}{4. Метод узловых точек}}

\key{Узловые точки} --- точки графика, координаты которых легко и точно считываются.

Главный принцип:

\[
\boxed{
\text{число независимых условий}
=
\text{число неизвестных коэффициентов}
}
\]

Например:

\[
f(x)=kx+b
\]
содержит два неизвестных коэффициента, поэтому обычно достаточно двух точек.

Для
\[
f(x)=2x^2+bx+c
\]
неизвестны $b$ и $c$, поэтому также достаточно двух точек.

\subsection*{Какие точки выбирать}

Лучше брать точки:

\begin{itemize}
  \item с целыми координатами;
  \item с координатами, близкими к нулю;
  \item с $x=0$ или $y=0$, если такие точки доступны;
  \item приводящие к простым вычислениям.
\end{itemize}

\sepLine

% ==================================================
\section*{\textcolor{darkaccent}{5. Универсальный алгоритм}}

\begin{enumerate}
  \item Определите тип функции.
  \item Запишите её общую формулу.
  \item Посчитайте неизвестные коэффициенты.
  \item Выберите нужное число узловых точек.
  \item Считайте координаты точек в порядке $(x;y)$.
  \item Подставьте координаты в формулу.
  \item Решите полученную систему.
  \item Соберите формулу функции.
  \item Выполните требуемое действие.
  \item Проверьте результат по рисунку.
\end{enumerate}

\sepLine

% ==================================================
\section*{\textcolor{darkaccent}{6. Пример: прямая}}

Пусть прямая проходит через точки
\[
A(-1;-3),
\qquad
B(3;4).
\]

Для
\[
f(x)=kx+b
\]
получаем
\[
\begin{cases}
-k+b=-3,\\
3k+b=4.
\end{cases}
\]

Отсюда
\[
k=\frac74,
\qquad
b=-\frac54.
\]

Значит,
\[
\boxed{f(x)=\frac74x-\frac54}.
\]

Например,
\[
\boxed{f(-5)=-10}.
\]

\subsection*{Быстрый способ найти наклон}

Если известны точки $(x_1;y_1)$ и $(x_2;y_2)$, то
\[
\boxed{k=\frac{y_2-y_1}{x_2-x_1}},
\qquad x_1\ne x_2.
\]

Также
\[
k=\tg\alpha.
\]

\sepLine

% ==================================================
\section*{\textcolor{darkaccent}{7. Пересечение графиков}}

В точке пересечения значения функций совпадают:

\[
\boxed{f(x)=g(x)}.
\]

Найдя $x$, подставьте его в любую функцию и получите $y$.

\sepLine

% ==================================================
\section*{\textcolor{darkaccent}{8. Парабола}}

Для
\[
f(x)=2x^2+bx+c
\]
точка с $x=0$ особенно удобна, потому что
\[
f(0)=c.
\]

Если график проходит через $(0;-4)$ и $(1;1)$, то
\[
c=-4,
\qquad
1=2+b-4,
\]
поэтому
\[
b=3.
\]

Следовательно,
\[
\boxed{f(x)=2x^2+3x-4}.
\]

\sepLine

% ==================================================
\section*{\textcolor{darkaccent}{9. Гипербола}}

Для функции
\[
f(x)=\frac{k}{x}+a
\]
обязательно
\[
\boxed{x\ne0}.
\]

Если график проходит через $(-3;0)$ и $(1;4)$, то
\[
\begin{cases}
0=-\dfrac{k}{3}+a,\\
4=k+a.
\end{cases}
\]

Отсюда
\[
a=1,
\qquad
k=3,
\]
поэтому
\[
\boxed{f(x)=\frac3x+1},
\qquad x\ne0.
\]

\sepLine

% ==================================================
\section*{\textcolor{darkaccent}{10. Частые ошибки}}

\begin{enumerate}
  \item Путают $x$ и $y$ в координатах точки.
  \item Путают $f(a)$ и $f(x)=a$.
  \item Теряют знак при переносе.
  \item Выбирают слишком неудобные узловые точки.
  \item Забывают ОДЗ у дробей, корней и логарифмов.
  \item Не проверяют полученный ответ по рисунку.
\end{enumerate}

% ==================================================
% ТРЕНИРОВОЧНЫЙ БЛОК
% ==================================================
\clearpage
\section*{\textcolor{darkaccent}{Тренировочный блок}}

Во всех задачах координаты узловых точек нужно считать непосредственно с графика.

\subsection*{Средний уровень}

\begin{enumerate}[label=\textbf{\arabic*)}]

% --------------------------------------------------
\item
На рисунке изображён график
\[
f(x)=kx+b.
\]
Восстановите формулу и найдите $f(6)$.

\begin{center}
\begin{tikzpicture}
\begin{axis}[
  schoolaxis,
  xmin=-4,xmax=4,
  ymin=-2,ymax=7,
  xtick distance=1,
  ytick distance=1
]
\addplot[maincurve,domain=-4:4]{x+3};
\addplot[knot] coordinates {(-2,1) (2,5)};
\node[above left] at (axis cs:-2,1) {$A$};
\node[above right] at (axis cs:2,5) {$B$};
\end{axis}
\end{tikzpicture}
\end{center}

% --------------------------------------------------
\item
На рисунке изображён график
\[
f(x)=kx+b.
\]
Восстановите формулу и найдите $f(0)$.

\begin{center}
\begin{tikzpicture}
\begin{axis}[
  schoolaxis,
  xmin=-5,xmax=3,
  ymin=-6,ymax=6,
  xtick distance=1,
  ytick distance=1
]
\addplot[maincurve,domain=-4:2]{-2*x-2};
\addplot[knot] coordinates {(-3,4) (1,-4)};
\node[above left] at (axis cs:-3,4) {$A$};
\node[below right] at (axis cs:1,-4) {$B$};
\end{axis}
\end{tikzpicture}
\end{center}

\clearpage

% --------------------------------------------------
\item
На рисунке изображена парабола
\[
f(x)=x^2+bx+c.
\]
Восстановите формулу и найдите $f(-1)$.

\begin{center}
\begin{tikzpicture}
\begin{axis}[
  schoolaxis,
  xmin=-4,xmax=4,
  ymin=-5,ymax=8,
  xtick distance=1,
  ytick distance=1
]
\addplot[maincurve,domain=-3.5:3]{x^2+x-3};
\addplot[knot] coordinates {(0,-3) (2,3)};
\node[below right] at (axis cs:0,-3) {$A$};
\node[above right] at (axis cs:2,3) {$B$};
\end{axis}
\end{tikzpicture}
\end{center}

\subsection*{Выше среднего}

% --------------------------------------------------
\item
На рисунке изображены две прямые $y=f(x)$ и $y=g(x)$.
Восстановите обе формулы и найдите абсциссу точки пересечения.

\begin{center}
\begin{tikzpicture}
\begin{axis}[
  schoolaxis,
  xmin=-3,xmax=4,
  ymin=-2,ymax=10,
  xtick distance=1,
  ytick distance=1
]
\addplot[maincurve,domain=-2.5:2.8]{2*x+4};
\addplot[secondcurve,domain=-1:4]{-2*x+7};
\addplot[knot] coordinates {(-1,2) (2,8) (0,7) (3,1)};
\node[below left] at (axis cs:-1,2) {$A$};
\node[above left] at (axis cs:2,8) {$B$};
\node[above right] at (axis cs:0,7) {$C$};
\node[below right] at (axis cs:3,1) {$D$};
\end{axis}
\end{tikzpicture}
\end{center}

\clearpage

% --------------------------------------------------
\item
На рисунке изображён график
\[
f(x)=\frac{k}{x}+a.
\]
Восстановите формулу и найдите $f(4)$.

\begin{center}
\begin{tikzpicture}
\begin{axis}[
  schoolaxis,
  xmin=-5,xmax=5,
  ymin=-6,ymax=8,
  xtick distance=1,
  ytick distance=1
]
\addplot[maincurve,domain=-5:-0.5]{4/x+1};
\addplot[maincurve,domain=0.5:5]{4/x+1};
\addplot[knot] coordinates {(-2,-1) (1,5)};
\node[below left] at (axis cs:-2,-1) {$A$};
\node[above right] at (axis cs:1,5) {$B$};
\end{axis}
\end{tikzpicture}
\end{center}

% --------------------------------------------------
\item
На рисунке изображена парабола
\[
f(x)=2x^2+bx+c.
\]
Восстановите формулу и найдите $f(3)$.

\begin{center}
\begin{tikzpicture}
\begin{axis}[
  schoolaxis,
  xmin=-4,xmax=2,
  ymin=-1,ymax=12,
  xtick distance=1,
  ytick distance=1
]
\addplot[maincurve,domain=-3.4:1.2]{2*x^2+5*x+5};
\addplot[knot] coordinates {(0,5) (-2,3)};
\node[above right] at (axis cs:0,5) {$A$};
\node[above left] at (axis cs:-2,3) {$B$};
\end{axis}
\end{tikzpicture}
\end{center}

\clearpage

% --------------------------------------------------
\item
На рисунке изображена прямая
\[
f(x)=kx+b.
\]
Восстановите формулу и найдите абсциссу точки пересечения с осью $Ox$.

\begin{center}
\begin{tikzpicture}
\begin{axis}[
  schoolaxis,
  xmin=-6,xmax=4,
  ymin=-8,ymax=6,
  xtick distance=1,
  ytick distance=1
]
\addplot[maincurve,domain=-5.5:3]{-1.5*x-3};
\addplot[knot] coordinates {(-4,3) (2,-6)};
\node[above left] at (axis cs:-4,3) {$A$};
\node[below right] at (axis cs:2,-6) {$B$};
\end{axis}
\end{tikzpicture}
\end{center}

% --------------------------------------------------
\item
На рисунке изображён график
\[
f(x)=\frac{k}{x}+a.
\]
Восстановите формулу и найдите $x$, если $f(x)=0,2$.

\begin{center}
\begin{tikzpicture}
\begin{axis}[
  schoolaxis,
  xmin=-8,xmax=6,
  ymin=-5,ymax=6,
  xtick distance=1,
  ytick distance=1
]
\addplot[maincurve,domain=-8:-0.6]{4/x+1};
\addplot[maincurve,domain=0.6:6]{4/x+1};
\addplot[knot] coordinates {(-4,0) (2,3)};
\node[below left] at (axis cs:-4,0) {$A$};
\node[above right] at (axis cs:2,3) {$B$};
\end{axis}
\end{tikzpicture}
\end{center}

\clearpage

% --------------------------------------------------
\item
На рисунке изображены графики
\[
f(x)=x^2+bx+c,
\qquad
g(x)=mx+n.
\]
Восстановите обе функции и найдите все абсциссы точек пересечения.

\begin{center}
\begin{tikzpicture}
\begin{axis}[
  schoolaxis,
  xmin=-5,xmax=3,
  ymin=-6,ymax=9,
  xtick distance=1,
  ytick distance=1
]
\addplot[maincurve,domain=-4.3:2]{x^2+3*x-2};
\addplot[secondcurve,domain=-3.5:3]{2*x+2};
\addplot[knot] coordinates {(0,-2) (1,2) (-1,0) (1,4)};
\node[below right] at (axis cs:0,-2) {$A$};
\node[below right] at (axis cs:1,2) {$B$};
\node[above left] at (axis cs:-1,0) {$C$};
\node[above right] at (axis cs:1,4) {$D$};
\end{axis}
\end{tikzpicture}
\end{center}

\subsection*{Сложный уровень}

% --------------------------------------------------
\item
На рисунке изображены
\[
f(x)=\frac{k}{x}+a,
\qquad
g(x)=mx+n.
\]

\begin{enumerate}[label=\arabic*)]
  \item восстановите обе формулы;
  \item найдите абсциссы точек пересечения;
  \item найдите координаты точек пересечения.
\end{enumerate}

\begin{center}
\begin{tikzpicture}
\begin{axis}[
  schoolaxis,
  xmin=-5,xmax=5,
  ymin=-5,ymax=6,
  xtick distance=1,
  ytick distance=1
]
\addplot[maincurve,domain=-5:-0.45]{2/x};
\addplot[maincurve,domain=0.45:5]{2/x};
\addplot[secondcurve,domain=-5:4.5]{x+1};
\addplot[knot] coordinates {(-1,-2) (2,1) (-2,-1) (2,3)};
\node[below left] at (axis cs:-1,-2) {$A$};
\node[below right] at (axis cs:2,1) {$B$};
\node[above left] at (axis cs:-2,-1) {$C$};
\node[above right] at (axis cs:2,3) {$D$};
\end{axis}
\end{tikzpicture}
\end{center}

\end{enumerate}

% ==================================================
% ОТВЕТЫ
% ==================================================
\clearpage
\section*{\textcolor{darkaccent}{Ответы}}

\begin{center}
\renewcommand{\arraystretch}{1.35}
\begin{tabular}{@{}c p{0.82\linewidth}@{}}
\toprule
\textbf{№} & \textbf{Ответ} \\
\midrule
1 & $f(x)=x+3$, \quad $f(6)=9$. \\
2 & $f(x)=-2x-2$, \quad $f(0)=-2$. \\
3 & $f(x)=x^2+x-3$, \quad $f(-1)=-3$. \\
4 & $f(x)=2x+4$, $g(x)=-2x+7$, \quad $x=\dfrac34$. \\
5 & $f(x)=\dfrac4x+1$, \quad $f(4)=2$. \\
6 & $f(x)=2x^2+5x+5$, \quad $f(3)=38$. \\
7 & $f(x)=-\dfrac32x-3$, \quad $x=-2$. \\
8 & $f(x)=\dfrac4x+1$, \quad $x=-5$. \\
9 & $f(x)=x^2+3x-2$, $g(x)=2x+2$, \quad $x_{1,2}=\dfrac{-1\pm\sqrt{17}}2$. \\
10 & $f(x)=\dfrac2x$, $g(x)=x+1$; точки пересечения: $(-2;-1)$ и $(1;2)$. \\
\bottomrule
\end{tabular}
\end{center}

\vfill

\begin{center}
\textcolor{darkaccent}{\textbf{Главная схема}}

\[
\boxed{
\text{график}
\longrightarrow
\text{узловые точки}
\longrightarrow
\text{система}
\longrightarrow
\text{формула}
\longrightarrow
\text{ответ}
}
\]
\end{center}

\end{document}
